\documentclass[12pt]{article}
\usepackage{sbc-template}
\usepackage[utf8]{inputenc}
\usepackage[T1]{fontenc}
\usepackage[brazil]{babel}
\usepackage{graphicx,url,amsmath}
\sloppy
\title{Mineração de métricas DORA em repositórios\\ populares com GitHub Actions}
\author{Víctor G. C. Pereira\inst{1}, Jonathan S. da Silva\inst{1},\\
Matheus F. de Oliveira\inst{1}}
\address{Laboratório de Medição e Experimentação de Software\\
Engenharia de Software -- PUC Minas}
\begin{document}
\maketitle
\begin{abstract}
This study defines a reproducible protocol to mine delivery metrics from public
GitHub repositories using GitHub Actions. Published releases represent deployments,
while workflow runs indicate failures and recoveries. The protocol uses a twelve-month
window, explicit inclusion criteria, and diagnostics for incomplete collections.
This draft presents the study structure and prior hypotheses; empirical results
and the complete DORA analysis remain subject to collection and validation.
\end{abstract}
\begin{resumo}
Este estudo define um protocolo reproduzível para minerar métricas de entrega em
repositórios públicos do GitHub que usam Actions. Releases publicadas representam
implantações, e runs indicam falhas e recuperações. O protocolo utiliza uma janela
de doze meses, critérios explícitos de inclusão e diagnósticos de coleta incompleta.
Este rascunho apresenta a estrutura do estudo e hipóteses prévias; os resultados
empíricos e a análise DORA completa dependem da coleta e validação.
\end{resumo}
% Seção 1 do artigo do Lab03 (template SBC).
% Hipóteses escritas antes da coleta dos dados: a data do commit que criou este
% arquivo no repositório registra esse momento.

\section{Introdução}
\label{sec:introducao}

A adoção de práticas DevOps tornou a entrega contínua de software um objetivo
comum em equipes de desenvolvimento. Para medir o desempenho dessa entrega, o
programa de pesquisa DORA (\textit{DevOps Research and Assessment}) propôs quatro
métricas: frequência de implantação (\textit{deployment frequency}), tempo de
espera para mudanças (\textit{lead time for changes}), taxa de falha de mudanças
(\textit{change failure rate}) e tempo de recuperação de serviço
\cite{forsgren2018accelerate}. A partir delas, os relatórios anuais do DORA
classificam as organizações em quatro grupos de desempenho: Elite, High, Medium
e Low \cite{dora2021}.

Essas classificações, porém, vêm de questionários respondidos por profissionais,
e não de dados medidos diretamente. Em projetos de código aberto, boa parte das
informações necessárias para calcular as métricas é pública: as releases marcam
as entregas, os commits mostram quando cada mudança foi feita, e as execuções de
integração contínua registram falhas e correções. O GitHub Actions, serviço de
automação de fluxos de trabalho integrado ao GitHub, tornou-se amplamente usado
nesses projetos \cite{kinsman2021actions, decan2022actions}, e o histórico de
suas execuções pode ser consultado pela API da plataforma.

Calcular as métricas DORA a partir desses dados exige decisões que os relatórios
não detalham. Em um projeto de código aberto não há implantação em produção
controlada pela equipe, então é preciso escolher o que conta como entrega, como
o lead time é medido e o que caracteriza uma falha. Cada escolha pode levar a
valores e classificações diferentes para o mesmo repositório.

O objetivo deste trabalho é calcular as quatro métricas DORA para repositórios
populares de código aberto que usam GitHub Actions, a partir de dados públicos
coletados pela API do GitHub, em uma janela de observação de 12 meses. Neste
estudo, uma implantação corresponde a uma release publicada, e as falhas são
observadas nas execuções de workflows do branch principal. Além de descrever a
distribuição das métricas e a classificação resultante, o estudo investiga a
relação entre as métricas, a influência de características dos repositórios e o
quanto a classificação depende das definições operacionais escolhidas.

\subsection{Questões de pesquisa e hipóteses}
\label{sec:hipoteses}

As quatro primeiras questões de pesquisa descrevem cada métrica na amostra. As
hipóteses abaixo são informais e foram registradas antes da coleta dos dados.
Elas servem como ponto de comparação na discussão dos resultados.

\textbf{RQ01 -- Qual é a frequência de implantação dos repositórios?}
Esperamos que a maioria dos repositórios fique na classe High, com mediana entre
uma e três releases por mês, e que pouquíssimos atinjam a classe Elite. Projetos
de código aberto agrupam mudanças em versões, e mesmo projetos muito ativos
raramente publicam uma release por dia. Também esperamos uma distribuição
assimétrica, com uma cauda longa de repositórios que publicam com muito mais
frequência do que a mediana.

\textbf{RQ02 -- Qual é o lead time das mudanças?}
Esperamos que a mediana do lead time por commit fique entre uma e três semanas,
o que coloca a maior parte dos repositórios na classe Medium. Se as releases
saem aproximadamente uma vez por mês, um commit espera, em média, metade desse
intervalo até ser publicado. Esperamos ainda que a variante calculada por release
produza valores maiores que a calculada por commit, porque ela é influenciada
pelos commits mais antigos de cada release.

\textbf{RQ03 -- Qual é a taxa de falha das mudanças?}
Pela variante baseada nas execuções de workflows, esperamos uma mediana entre
10\% e 20\%. Falhas de integração contínua são frequentes em projetos ativos e
nem sempre indicam um problema que chegaria ao usuário, como testes instáveis ou
erros de configuração. Esperamos que a variante baseada em releases corretivas
produza valores diferentes, e que as duas variantes classifiquem o mesmo
repositório de forma diferente em uma parcela relevante dos casos.

\textbf{RQ04 -- Qual é o tempo de recuperação após uma falha?}
Esperamos uma mediana de algumas horas, abaixo de um dia, o que coloca a maior
parte dos repositórios na classe High. Em projetos ativos, uma execução com falha
costuma ser seguida por uma correção em pouco tempo, já que ela bloqueia o
trabalho dos demais contribuidores. Esperamos também uma cauda longa, formada
por episódios que só se resolvem dias depois, e alguns episódios ainda abertos
no fim da janela.

\section{Metodologia}
\label{sec:metodologia}
\subsection{Seleção e coleta}
A população de candidatos é obtida pela Search API do GitHub, com pelo menos
mil estrelas, sem forks nem repositórios arquivados e com atividade desde o
início da janela. A janela é o intervalo UTC de 1 de outubro de 2025 (inclusivo)
a 1 de outubro de 2026 (exclusivo). Incluímos projetos com workflows próprios
em \texttt{.github/workflows/}, pelo menos cinco releases estáveis publicadas e cinquenta
runs válidos dentro da janela. Pré-releases são preservadas para variantes, mas ficam fora da definição
principal de deploy e inclusão; drafts são excluídos.

A coleta registra metadados, releases e workflow runs do default branch,
com evento \texttt{push}. As consultas de runs são mensais; a integração
subdivide intervalos saturados para lidar com o teto de mil resultados por
consulta. Os IDs são deduplicados. Cache por página e checkpoints permitem
retomada, e a incompletude é registrada, sem ser convertida em ausência de falhas.
O comando integrado implementa a seleção dos primeiros cem projetos elegíveis
completos na ordem da busca, que prioriza estrelas. Na execução disponível,
contudo, usamos o piloto já coletado: a coleta adicional foi encerrada por
limite de tempo. Cem candidatos não equivalem a cem incluídos, e a meta de
cem elegíveis completos não foi atingida.

\subsection{Definições operacionais e análise}
Deployment frequency é o número de releases estáveis publicadas dividido
pelas semanas reais da janela, calculadas a partir de sua duração em segundos. A CFR (a) é calculada pela Equação~\ref{eq:cfr}; são falhas as conclusions
\texttt{failure}, \texttt{timed\_out} e \texttt{startup\_failure}, e sucesso a
conclusion \texttt{success}. As demais conclusions são ignoradas. Sem runs
válidos, a CFR fica ausente.
\begin{equation}
\mathrm{CFR}(a)=\frac{\text{falhas}}{\text{falhas}+\text{sucessos}}
\label{eq:cfr}
\end{equation}
Os runs criados na janela são ordenados por \texttt{run\_started\_at} e ID em
cada workflow. A primeira falha após um sucesso observado inicia o episódio;
a duração é \texttt{updated\_at} do próximo sucesso menos
\texttt{run\_started\_at} da primeira falha. Falhas consecutivas mantêm o início.
Sem término antes do fim da janela, o episódio tem censura à direita, e
reportamos sua proporção por repositório. Sequências iniciais sem sucesso
anterior observado têm censura à esquerda separada, sem duração presumida.
Timestamps essenciais inválidos deixam o workflow sem estimativa, com
diagnóstico, evitando unir episódios ao remover somente um run.
A mediana e o intervalo interquartil (IQR), em horas, usam apenas recuperações
observadas, com interpolação linear inclusiva; não são estimativas de sobrevivência.

A coleta de commits compara releases estáveis consecutivas, recuperando a
release anterior à janela quando existir. O protocolo usa a data do autor do
commit. A consulta é paginada e registra comparações inacessíveis ou incompletas;
a primeira release histórica sem anterior é ignorada. O lead time por release
é a publicação menos a data do autor mais antiga entre os commits da comparação.
O valor do repositório é a mediana por release, sem ponderação por número de
commits. Ausência de commits, datas inválidas, comparações incompletas e tempos
negativos não produzem valor; zero observado é válido. Essa variante está
implementada e aguarda validação real. A variante por commit ainda depende de
integração e validação antes da análise conjunta do protocolo.
A classificação usa os cortes fixos do enunciado da disciplina: frequência
Elite a partir de sete releases/semana, High a partir de uma/semana e Medium
a partir de uma/mês; para lead time, os limites são 24, 168 e 720 horas.
CFR usa 15\%, 30\% e 45\%, e recuperação usa 1, 24 e 168 horas.
Os rótulos de classe originalmente previstos nas hipóteses ficam preservados
como registro anterior, mesmo onde divergiam dos cortes agora confirmados.

\subsection{Artefatos e rastreabilidade}
\begingroup\raggedright
O código, as decisões e os testes estão no repositório público:
\url{https://github.com/Victorgabrielcruz/lab-medicao-e-experimentacao-de-software}.\par
O planejamento pode ser consultado no GitHub Projects:
\url{https://github.com/users/Victorgabrielcruz/projects/7}.\par
O relatório técnico está em \texttt{Lab03/docs/relatorio-victor-sprint-01.md}.\par
Dados brutos, cache, credenciais e saídas reais permanecem locais. Hashes da
entrada e contagens agregadas permitem identificar a execução auditada.
\par\endgroup

\section{Resultados disponíveis do piloto}
\label{sec:resultados}
O piloto partiu de cem candidatos, com 82 repositórios após a verificação
de Actions e metadados. O snapshot consolidado contém 238.629 runs
deduplicados no recorte. A coleta de runs é completa em 64 repositórios e
incompleta em 18, condição que permanece explícita nas saídas.

Para a CFR (a), há 217.412 runs válidos: 22.893 falhas e 194.519 sucessos;
21.217 runs têm conclusions ignoradas. Há valor observado em 70 repositórios
e ausência em doze. Essas contagens agregadas não substituem a distribuição
da CFR por repositório nem permitem generalizar para a população.

Na RQ04 corrigida conforme o enunciado, foram auditados 6.163 episódios
recuperados e 38 censurados à direita, em um total de 6.201 episódios definidos.
Há recuperação observada em 56 repositórios. Registram-se também 92 sequências
iniciais com censura à esquerda e um repositório com dados temporais
incompletos. Os históricos parciais podem alterar as contagens e os episódios;
a auditoria demonstra correspondência com a entrada, sem provar completude.

Uma leitura exclusivamente do cache existente, sem novas consultas,
recuperou releases com paginação completa para cinco repositórios; os demais
77 não têm essa evidência disponível. Combinando os dados de runs e releases,
nenhum repositório pôde ser confirmado como elegível completo, 42 já não
atendem aos critérios conhecidos e 40 têm elegibilidade indeterminada.
Esses diagnósticos não são etapas cumulativas de um funil.
Não houve observação suficiente para lead time por release.

A execução adicional destinada a cem elegíveis foi encerrada por limite de
tempo, preservando o material já obtido. Assim, esta entrega é parcial e
não atende ao critério de cem projetos elegíveis completos. Os testes com
fixtures verificam propriedades da implementação, sem constituir resultados
empíricos. A confirmação ou refutação das hipóteses e as análises conjuntas
das RQs permanecem pendentes.

\section{Discussão e ameaças à validade}
\label{sec:discussao}
\subsection{Validade de construto}
Uma release pública é uma aproximação de implantação, e falhas de workflows
não identificam necessariamente incidentes em produção. Runs podem incluir
testes instáveis, falhas de infraestrutura ou de configuração. Recuperação
observada em CI não equivale à restauração de serviço percebida pelo usuário.
As variantes operacionais e as classes fixadas pelo enunciado devem ser comparadas somente
após integração e validação das métricas correspondentes.

\subsection{Validade interna e confiabilidade}
Limites de paginação, indisponibilidade de repositórios e retenção ou remoção
de históricos podem produzir dados incompletos mesmo numa janela fechada.
Releases podem ser editadas ou publicadas posteriormente; uma janela fechada
não garante imutabilidade dos registros. Cache preserva o snapshot observado,
mas deve ser distinguido de uma consulta atual. O default branch é observado
na coleta e pode ter sido diferente no passado. A censura e recuperações
anteriores ao início da janela exigem cuidado ao interpretar as durações.

\subsection{Validade externa}
Repositórios populares, usuários de Actions e projetos com releases frequentes
não representam todos os projetos de software. Selecionar os primeiros elegíveis
na ordem de estrelas introduz viés para projetos populares. As conclusions devem
ser limitadas à população e à janela estudadas. A discussão empírica das hipóteses
será adicionada após a auditoria, sem alterar seu registro anterior à coleta.

\subsection{Validade de conclusão}
A análise futura deve tratar comparações múltiplas, efeitos de subgrupos pequenos
e durações censuradas. As hipóteses complementares foram registradas depois do
piloto, o que limita sua interpretação como previsões anteriores aos dados.
Nenhum teste estatístico ou veredito das hipóteses é apresentado nesta sprint.

\section{Conclusão}
\label{sec:conclusao}
O estudo estabelece decisões operacionais e uma estrutura reproduzível para
coletar e calcular métricas de entrega em projetos que usam GitHub Actions.
A conclusão empírica depende da amostra elegível completa e da integração das
variantes ainda pendentes. Após a validação, esta seção deve retomar cada RQ,
sintetizar os resultados observados e delimitar as implicações das aproximações
adotadas. Nesta versão, não se afirma o atendimento das hipóteses.

\bibliographystyle{sbc}
\bibliography{referencias}
\end{document}
